\newcommand{\figuraCierreAutogravitante}{%
\begin{figure}[H]
    \centering
    \begin{tikzpicture}[>=Latex,font=\small]
        \node[font=\bfseries] at (6.25,5.45)
            {Acoplamiento de un medio continuo autogravitante};

        \node[draw=blue!65!black,very thick,rounded corners=3pt,
              fill=blue!7,minimum width=3.0cm,minimum height=1.05cm,
              align=center] (rho) at (1.85,3.55)
            {distribución\\de densidad};
        \node[draw=orange!75!black,very thick,rounded corners=3pt,
              fill=orange!8,minimum width=3.4cm,minimum height=1.05cm,
              align=center] (phi) at (6.20,3.55)
            {potencial\\gravitacional};
        \node[draw=green!45!black,very thick,rounded corners=3pt,
              fill=green!8,minimum width=3.3cm,minimum height=1.05cm,
              align=center] (press) at (10.65,3.55)
            {equilibrio\\hidrostático};
        \node[draw=red!70!black,very thick,rounded corners=3pt,
              fill=red!6,minimum width=3.5cm,minimum height=1.05cm,
              align=center] (eos) at (6.20,1.15)
            {ecuación\\de estado};

        \draw[->,ultra thick,blue!60!black]
            (rho)--(phi)
            node[midway,above] {fuente};
        \draw[->,ultra thick,orange!75!black]
            (phi)--(press);
        \draw[->,ultra thick,green!45!black]
            (press) to[out=-120,in=15] (eos);
        \draw[->,ultra thick,red!70!black]
            (eos) to[out=165,in=-60] (rho);
    \end{tikzpicture}
    \caption{La densidad genera el potencial gravitacional mediante la
    ecuación de Poisson; el potencial determina el equilibrio de presión, y
    la ecuación de estado cierra el sistema relacionando nuevamente presión y
    densidad.}
    \label{fig:cierre-autogravitante}
\end{figure}%
}